\documentclass[../main.tex]{subfiles}
\begin{document}
%==========================================TIÊU ĐỀ TẬP========================================
\begin{table}[h]
  \begin{adjustwidth}{-1cm}{}
    \begin{tabular}{>{\centering\arraybackslash}p{6cm}|>{\raggedright\arraybackslash}p{12.5cm}}
      \multirow{5}{6cm}
      % Cột bên trái: ÉPISODE và số tập
      {\\[-40pt]
        \flaregothic\fontsize{20pt}{20pt}\selectfont \centering
        {\contour{errortwo}{\textcolor{black}{ÉPISODE}}}\color{black}\\
        \burbank\fontsize{72pt}{72pt}\selectfont
        \includegraphics[height=3cm]{../../../../Image/Book?/number/num1.png}
        \\[18pt]
        % \flaregothic\fontsize{14pt}{14pt}\selectfont
        % \color{epattr}BIENVENUE, \uline{\color{Banpire} Mel} !
        % \flaregothic\fontsize{18pt}{18pt}\selectfont
        % \color{epattr}ÉPISODE PERDU
        \color{black}
      }
       &
      % Dòng thứ nhất: tên tập tiếng Nhật
      {\fotskip\fontsize{18pt}{18pt}\selectfont \ruby{青}{あお}\ruby{上}{がみ}\ruby{高}{こう}\ruby{校}{こう}へようこそ
      } \\[6pt]
      % Dòng thứ hai: tên tập tiếng Anh
       & {\cafeteria\fontsize{18pt}{18pt}\selectfont Welcome to Aogami High!}                        \\[4pt]
      % Dòng thứ ba: tên tập tiếng Việt
       & {\vnmsans\fontsize{18pt}{18pt}\selectfont Chào mừng tới Trường Trung học Aogami!}           \\[8pt]
      % Dòng thứ tư: Nhân vật xuất hiện
       & {\caviar{\fontsize{14pt}{14pt}\selectfont Track used: Dream of Sky - Łukasz Michalski}}
      \\[6pt]
      % Dòng cuối cùng: Thời gian diễn ra của tập (thường sẽ là ngày phát hành)
       & {\continuum\fontsize{16pt}{16pt}\selectfont Ngày phát hành: 18/02/2022}                     \\[8pt]
    \end{tabular}
  \end{adjustwidth}
\end{table}
%===========================================NỘI DUNG==========================================
\color{black}

\begin{center}
  \pov{Haragen Kyuen}
  \uline{Nhà Haragen.\\Năm giờ sáng.}
\end{center}

\dots

``\dots Ư\dots''

\href{https://www.youtube.com/watch?v=Mqokmn0WOMo}{Tiếng chuông báo thức của điện thoại tôi kêu.}

Nhiều người sẽ gọi đây là thứ âm thanh đáng ghét, ám ảnh, thậm chí là ``quỷ dữ'' mỗi sáng. Nhưng với tôi\dots\ nó lại khá là nhẹ nhàng, thậm chí có chút êm ái. Hoặc cũng có thể\dots\ chỉ mình tôi nghĩ vậy thôi.

\dots

\dots

\dots Được rồi, thú thật thì cái báo thức kia cũng khiến tôi hơi khó chịu thật. Nhất là khi ngày nào cũng phải bò dậy vào cái giờ oái oăm thế này.

``Ừm\dots'' tôi rên khẽ, quờ tay tìm chiếc điện thoại vẫn đang rung bần bật đâu đó trên đầu giường.

``\dots\ Ầy da\dots\ Ay! Đây rồi!''

Tôi hé mắt ra, ánh sáng từ màn hình lập tức đập vào, báo năm giờ sáng. Đúng là chẳng ai tầm tuổi tôi lại dậy vào giờ này, nhưng không hiểu sao tôi đã quen rồi.

Bên cạnh tôi, Saikyou - em gái tôi - vẫn cuộn tròn trong chăn, mái tóc xoăn bung xõa trên gối, nhịp thở đều đặn. Không giống tôi, con bé thuộc loại ``ngủ được là ngủ tới trưa''.

Tôi chậm rãi ngồi dậy, dựa lưng vào thành giường, tắt chuông báo thức và khẽ thở dài. Ngoài kia vẫn tối om, mặt trời chưa ló dạng. Chắc các bạn đang tự hỏi: tại sao một thằng nhóc mười sáu tuổi như tôi lại dậy sớm đến vậy?

Đơn giản thôi: hôm nay là ngày đầu tiên chúng tôi đến trường mới.

Thực ra, đây chẳng phải lần đầu chúng tôi phải làm quen với môi trường mới. Bố mẹ tôi hay bị thuyên chuyển công tác (có lẽ vì cấp trên bắt, tôi cũng không chắc), thành ra chuyện chuyển trường với chúng tôi là cơm bữa. Đếm sơ sơ chắc cũng đủ hai bàn tay, hoặc hơn. Tôi không nhớ rõ nữa\dots\ chắc vì tôi cũng chẳng để tâm lắm.

Lần này, chúng tôi chuyển tới một thị trấn miền núi tên là Aogami. Hôm nay sẽ là lần đầu bước chân vào Trường Trung học Aogami - ngôi trường mang cùng tên thị trấn.

Nghe đồn đây là một ngôi trường khá nổi tiếng ở vùng này. Không chỉ vì thành tích học tập hay tỷ lệ dỗ đại học cao hơn hẳn so với các trường khác, mà còn bởi những truyền thuyết kỳ dị mà dân địa phương truyền tai nhau.

Người ta bảo rằng\dots\ ở thị trấn này  có một lời nguyền.

Mà nhắc tới ``lời nguyền'', tôi lại nhớ tới một giấc mơ kỳ lạ mà tôi vừa trải qua đêm qua.

Trong mơ, tôi và Saikyou bị lạc vào một nơi trông rất giống trường học, nơi chúng tôi phải tìm cách thoát ra. Rồi một trận động đất ập tới, cuốn cả hai tới một chiều không gian khác, cũng là trường học, nhưng mang dáng dấp của quá khứ xa xôi, cách đây khoảng sáu, bảy thập kỷ về trước.

Chúng tôi gặp đủ loại chuyện quái đản: những bóng dáng học sinh tôi chưa từng gặp nhưng lại thấy quen đến rợn người; những vụ sập đổ bất thình lình; và\dots

\dots một cô gái. Thực ra là hai, nhưng tôi vẫn chưa chắc có nên gộp họ làm một hay không.

Cô gái tóc nâu ấy\dots\ ``Misora Shino'', nếu trí nhớ tôi không nhầm. Vóc dáng thanh mảnh, đôi mắt xanh trong vắt, và một nụ cười tưởng như thân thiện\dots\ nhưng tại sao cô ấy lại xuất hiện ở đó?

Và rồi, tôi lại gặp thêm một cô gái khác, gần như bản sao của Shino, chỉ khác đôi mắt màu nâu và bộ đồng phục khác hẳn. Cô tự xưng là Sora. Ước nguyện của cô thật giản dị: muốn được tận hưởng tuổi học trò, muốn có bạn, muốn được chơi đùa và cười đùa cùng ai đó. Giọng nói ấy, ánh mắt ấy\dots\ khiến tôi khó lòng từ chối.

Cuối cùng, chúng tôi đã cùng nhau bước qua một cánh cửa sáng rực, tin rằng mình đã tìm thấy lối ra\dots

\dots Nhưng rồi tôi mở mắt, và nhận ra tất cả chỉ là một giấc mơ.

Một giấc mơ kỳ quặc. Nếu không muốn nói là\dots\ hãi hùng.

Tôi vươn vai, định đứng dậy thì chợt nhận ra\dots\ trong túi quần ngủ có gì đó cộm lên. Kỳ lạ thật, tối qua tôi nhớ rõ mình đâu có bỏ gì vào đây.

Tò mò, tôi thò tay vào và lôi ra hai thứ.

Một sợi dây đeo đồng hồ màu đen đã bắt đầu hơi nhăn lại, và một mặt đồng hồ tròn trịa, hơi bụi.

Tôi khựng lại. Đây chẳng phải là\dots?

\dots Không, chắc chỉ là trùng hợp. Hoặc\dots\ tôi vẫn chưa tỉnh hẳn?

\dots

\dots

``Mình đang nghĩ gì chứ? Tại sao lại nghĩ tới mấy thứ kỳ quái đó vào ngày trọng đại này?'' tôi tự nhủ, cố gắng ngắt mạch suy nghĩ đang xoáy vào những hình ảnh trong mơ.

Đúng rồi, tôi còn phải đến trường, có lẽ sẽ gắn bó ở đây một năm. Bạn mới, giáo viên mới, môi trường mới\dots\ giống hệt những lần chuyển trường trước.

Sự thật là hai anh em tôi hầu như không có bạn. Nói chuyện xã giao thì có, nhưng để thân thiết thì\dots\ không. Một phần vì tôi hướng nội, phần lớn vì tính bảo vệ thái quá của em gái tôi. Nhiều lần tôi bị bắt nạt, và cũng bấy nhiêu lần em ấy nhảy vào mắng chửi, thậm chí đánh nhau với bọn kia.

Chẳng trách mẹ hay đùa: ``Bạn của các con chắc chỉ có mỗi các con thôi.'' Ngẫm nghĩ lại\dots\ bà ấy nói cũng có phần đúng.

Tôi liếc điện thoại: 5:15. Còn nhiều thời gian. Ở đây vào học từ 6:45, khác hẳn Nhật Bản. Tôi thì quen dậy sớm, nhưng Saikyou thì không.

``\dots Saikyou vẫn còn đang nướng cơ à,'' tôi lẩm bẩm, lay vai cô em gái. ``Nào, dậy đi học thôi.''

``Anh gọi sớm thế\dots'' em ấy rên rỉ, rồi ngồi dậy, ngáp dài. Tôi chống hông: ``Em sẽ phải quen thôi. Ở đây không giống như nơi ở cũ đâu.''

Cô bé lơ mơ nhìn quanh, như chưa tin mình đang ở thực tại. Tôi chẹp miệng: ``Ngốc ạ, đây là thực tại, không còn ở trong mơ đâu.''

Im lặng.

Rồi, giọng em ấy khẽ vang lên: ``\dots Em vừa mơ thấy một giấc mơ kỳ lạ, Nii-chan ạ.''

Tôi ngẩng lên. ``Là sao?''

``Em mơ thấy\dots\ hai chúng ta bị lạc vào một ngôi trường. Mắc kẹt. Phải tìm lối ra. Rồi rơi xuống một cái hố\dots\ quay về quá khứ\dots\ và mọi thứ đổ sập\dots''

Từng chữ như đâm thẳng vào ký ức tôi. Đó chính là giấc mơ của tôi.

``Anh cũng mơ thấy như vậy.''

Saikyou lập tức tỉnh hẳn. ``Thế á? Có phải vì ngủ chung giường không?''

Tôi đổ mồ hôi: ``Không nghĩ là vậy\dots''

``Nhưng\dots\ cảm giác nó thật lắm. Như mình đã ở đó thật.''

Cả hai nhìn nhau, đồng thời nhớ lại những cảnh tượng trong mơ.

``Em nhớ là chúng ta có nhìn thoáng qua những học sinh trong ngôi trường ấy\dots\ ở cả hiện đại và ở quá khứ. Họ cứ thoắt ẩn thoắt hiện\dots'' em bắt đầu ngồi kể lại những hiện tượng siêu nhiên mà chúng tôi đã bắt gặp. ``Thật sự luôn, em đã hoảng hốt vì anh bỗng chốc bị nhức đầu á!''

``\dots Đúng thật. Rồi sau đó là\dots\ hoa bỉ ngạn đỏ?''

``Có, em cũng thấy. Hai lần.''

``Và bức tranh khiến cho em ôm đầu mà hét lên một cách điên dại\dots''

Saikyou run rẩy khắp người, sợ đến xanh mặt. ``Bây giờ nhớ lại thời điểm đó thôi cũng khiếp em lẩy bẩy tay chân\dots''

``Và\dots''

``Hai cô gái kia,'' chúng tôi đồng thanh. Một là Shino. Một là Sora - giống hệt, chỉ khác màu mắt và bộ đồng phục. Sự xuất hiện của hai người đó khiến hai anh em chúng tôi xuất hiện đầy rẫy những thắc mắc.

Trong khoảnh khắc ấy, tôi chợt nhận ra mình vẫn đang cầm hai món đồ trong tay: dây đeo và mặt đồng hồ. Chúng lạnh ngắt, nặng trĩu, cứ như muốn chứng minh rằng\dots\ đó không chỉ là mơ.

Tại sao những điều đó lại xảy ra?

Chuyện quái gì đang xảy ra ở đây vậy?

Và những hiện tượng đó xảy ra mang điềm báo gì?

\dots

\dots

Thực sự, toàn những câu hỏi dạng ``tu từ'' như thế, không một ai có thể trả lời được ngay cả.

Ngoại trừ một ai đó. Một ai đó cụ thể biết hết sự vụ về việc này.

Nhưng đó là ai? Và người đó đang ở đâu?

Đó cũng là những câu hỏi hóc búa nữa mà chưa tìm được một lời giải đáp.

\dots

Mà, nghĩ làm gì cho mất công nữa. Tôi bỏ cuộc nghĩ về mấy thứ xa xăm như vậy. Giờ còn quan trọng hơn: chúng tôi phải chuẩn bị cho buổi học đầu tiên.

``Anh thấy ta chẳng cần nghĩ về mấy thứ đó nữa đâu. Thôi nào, ra khỏi giường chuẩn bị đồ đạc đi kẻo muộn.''

``À, phải. Hôm nay là ngày đầu tiên tới trường\dots'' Em gái tôi lồm cồm bò dậy, vừa nói vừa ngáp dài, rồi lẩm bẩm một tiếng ``lần nữa'' đầy chán chường.

``Chúng ta sẽ chuyển trường tới khi nào đây\~{}\dots''

``Chắc là cho đến khi ra trường quá. Thôi, không than vãn nữa.''

``Vâng\~{}\dots'' Saikyou đáp, giọng kéo dài như đang đọc thơ, rồi lười biếng lấy bộ đồng phục treo trên mắc và đi vào phòng tắm.

Tôi cũng đã chuẩn bị sẵn bộ đồng phục của mình từ tối hôm qua. Đồng phục của trường Aogami có hai loại: một cho nam, một cho nữ. Bộ của nam thì khá đơn giản, gồm áo sơ mi ngắn tay có túi áo ở phía trước ngực bên trái, quần dài xanh nước biển và cà vạt đồng màu. Bộ của nữ cũng gần tương tự, chỉ thay quần bằng váy ca-rô xanh (pha xanh thẫm) và thay cà vạt bằng chiếc nơ trước ngực. Cô em gái của tôi thậm chí còn cặp thêm một cái nơ đỏ trên tóc, trông vừa gọn gàng vừa\dots\ hơi đáng yêu, dù chắc em ấy sẽ không bao giờ thừa nhận.

Tiếng nước từ phòng tắm vọng ra không lâu, Saikyou bước ra với bộ đồng phục mới tinh, tóc vẫn còn ẩm. Khác với mấy cô gái mà tôi biết - những người có thể mất cả tiếng đồng hồ để chỉnh sửa tóc tai - em tôi lại cực kỳ nhanh gọn. Có lẽ vì đã quen với việc sáng ra là phải chạy đua với đồng hồ.

Tôi nhìn đồng hồ điện thoại, 6 giờ kém 5. Còn tận gần một tiếng nữa mới đến giờ vào lớp.

``Sách vở có đầy đủ hết chưa?'' tôi hỏi, vừa đeo cặp lên vai.

``Xong hết rồi\~{}'' Saikyou đáp, giọng tươi như thể chuyện chuyển trường năm lần bảy lượt chưa từng xảy ra.

``Nếu mà quên bút thì đừng có mà khóc lóc xin anh đâu đấy,'' tôi đùa.

``B-Biết rồi chứ!'' em ấy phồng má lên phản đối, nhưng tôi biết chắc kiểu gì đến lúc đó cũng sẽ có màn ``anh ơi cứu em''.

Tôi bật cười, mở cửa phòng ngủ và bước xuống tầng. Tiếng bước chân Saikyou lóc cóc theo sau, kèm theo tiếng thở dài của một kẻ chưa sẵn sàng chào đón buổi sáng.

\ct

Bữa sáng của chúng tôi cũng chẳng cầu kỳ gì, chỉ vài nắm cơm gói từ tối qua. Kể ra chúng tôi cũng quen rồi, bố mẹ làm việc xa, cả tháng mới ghé qua một, hai lần, nên chuyện hai anh em tự lo cho mình là điều hiển nhiên. Ban đầu cũng lúng túng, nhưng rồi tự lập trở thành thói quen\dots\ và cũng là một thứ ràng buộc đặc biệt giữa tôi và Saikyou.

Hai anh em chúng tôi cũng không có gì để ca thán cả. Trái lại, chúng tôi còn thấy vui là đằng khác. Được tự do làm chủ trong ngôi nhà này, ai mà chả muốn?

Dù đi đâu hay làm gì, tôi và em ấy luôn có nhau. Chúng tôi luôn làm việc cùng nhau, từ những việc nhỏ nhặt như nấu cơm, rửa bát, cho đến những việc lớn hơn một chút như dọn dẹp toàn bộ nhà cửa. Bọn tôi không muốn biến nhà mình thành một cái chuồng lợn, đúng chứ?

Và, kể cả là cùng nhau đi học. Có thể chúng tôi không có bạn, nhưng ít nhất chúng tôi vẫn còn có nhau. Vừa đi, vừa gặm cơm nắm, vừa có thể nhìn quanh khung cảnh thị trấn này, quả đúng là điều tuyệt vời.

Nhắc đến quang cảnh, thị trấn núi này không tấp nập như thành phố, nhưng lại thoáng đãng và sáng sủa. Mặt trời mới lên, ánh nắng len qua tán cây, mùi gỗ ẩm lẫn trong gió mát. Dù đã ở đây một thời gian, mỗi lần đi qua những con phố này, tôi vẫn có cảm giác như lần đầu.

\img{e1-1a}

Aogami là một thị trấn nhỏ nằm giữa những dãy núi xanh thẳm, nơi những ngôi nhà truyền thống bằng gỗ vẫn đứng sát cạnh các tòa nhà bê tông. Những con đường lát đá uốn lượn dọc theo sườn núi, dẫn tới những cửa hàng tạp hóa nhỏ xinh và những quán cà phê ấm cúng. Dù không có những trung tâm thương mại hay rạp chiếu phim hoành tráng, nhưng nơi đây vẫn có đầy đủ tiện nghi cơ bản của cuộc sống hiện đại.

Điều đặc biệt nhất của Aogami có lẽ là cách người dân nơi đây vẫn giữ gìn được nét văn hóa truyền thống giữa dòng chảy của thời đại. Những lễ hội theo mùa vẫn được tổ chức đều đặn, tiếng trống taiko vang vọng từ đền thờ cổ kính trên đỉnh đồi, trong khi những chiếc xe điện chạy êm ru trên những con phố dưới chân núi. Sự hòa quyện giữa cũ và mới tạo nên một bầu không khí thanh bình, khiến bất cứ ai đến đây đều cảm thấy như được trở về nhà.

Càng gần trường, càng đông học sinh nhập vào dòng người. Tiếng bước chân, tiếng trò chuyện ríu rít. Saikyou vừa nhai cơm vừa nhận xét: ``Hmm\dots\ nhiều người học ở đây ghê ha, Nii-chan.''

``Ừ. Hy vọng anh em mình sẽ kiếm được vài người bạn\dots''

``Em cũng hy vọng thế.''

Tất nhiên, tôi không quên chọc ghẹo: ``Nếu em bỏ được cái kiểu du côn đó thì hy vọng sẽ cao hơn nhiều.''

``B-Bởi vì em yêu Nii-chan quá nên mới thế chứ bộ!'' em ấy huơ tay phản đối, nhưng tôi chỉ biết thở dài.

Chúng tôi vừa cười vừa nói cho đến khi cổng trường Aogami hiện ra. Ba dãy nhà ba tầng đứng sừng sững, không mới nhưng cũng không cũ kỹ như tôi từng tưởng. Đây sẽ là nơi chúng tôi học, kết bạn, và tạo nên những kỷ niệm mới.

Ít nhất là, đó là điều tôi nghĩ\dots\ cho đến khi thấy \emph{cô ấy}.

\img{e1-1b}

Một cô gái tóc nâu, tết hai bím thấp, đứng yên bên lề cổng trường, như thể bị đông cứng giữa dòng người. Trên mái tóc ấy là một chiếc bờm tai chó cụp xuống. Cổ cô quấn một vòng dây đỏ, giống vòng cổ thú nuôi hơn là phụ kiện thời trang. Mắt nâu của cô vô hồn, miệng khẽ mấp máy.

Hình như chúng tôi đã gặp cô ấy ở đâu đó rồi thì phải.

``Mà, Nii-chan? Sao em thấy cô ấy trông quen quen nhỉ?''

``\dots Chắc thế. Lúc chúng ta mới chuyển vào đây đúng chứ?''

Hai chúng tôi nhìn cô nàng ``tai chó'' ấy, không hiểu tại sao lại đứng đó. Có thể cô ấy đang chờ người bạn của mình.

Tôi tiến tới, định chào hỏi cô ấy, nhưng rồi\dots

``\dots Mau trả bạn cho tôi, mau trả bạn cho tôi, mau trả bạn cho tôi\dots'' câu nói lặp đi lặp lại như một khúc nhạc lỗi.

Tôi và Saikyou thoáng khựng. Tôi nhớ mang máng từng gặp cô ấy khi mới chuyển đến đây, khi đó, cô tươi cười và thân thiện lắm.

Giờ thì\dots\ như thể đó là một người khác. Vẻ mặt của cô ấy trông không khác gì như đang gặp phải một cú sốc vậy.

``Nii-chan\dots\ em nghĩ là mình không nên\dots'' Saikyou khẽ kéo tay tôi.

Tôi gật nhẹ. Bước đi tiếp. Nhưng trong khoảnh khắc, tôi có cảm giác mình vừa lướt qua một thứ gì đó\dots\ như cái bóng lành lạnh trườn xuống sống lưng.

Cơ mà\dots\ không, hôm nay là buổi học đầu tiên. Tôi phải để ý vào chuyện khác thôi.

\ct

\begin{center}
  \uline{Hành lang tầng 2, dãy B - Trường Trung học Aogami.}
\end{center}

Hai anh em chúng tôi tiếp tục rảo bước vào lớp học. Tôi sẽ học lớp 1-C, và\dots

``\textbf{Nii-chan! Em lại được học chung với anh này! Nhìn đi!}''

\dots em gái tôi cũng vậy. Lần nữa. Đây không phải là lần đầu hai anh em tôi được học chung lớp. Thực tế thì, cứ mỗi lần chuyển trường, danh sách lớp lại ``vô tình'' sắp chúng tôi vào cùng nhau. Tôi còn chẳng biết đây là sự trùng hợp hay\dots\ một cái gì đó khác đang sắp đặt.

Và như mọi khi, ghế của chúng tôi lại sát nhau. Tôi đã quen với việc này đến mức không còn gì để phàn nàn nữa. Miễn là Saikyou thấy vui, thì tôi cũng chẳng bận tâm.

Chúng tôi rảo bước dọc hành lang, vừa tìm lớp vừa ngó nghiêng xung quanh. Khối năm nhất nằm ở tầng hai, tôi đã xem qua sơ đồ trường trước khi vào.

Saikyou thì hớn hở như thể đây là cuộc phiêu lưu mới, còn tôi thì\dots\ không hiểu sao lại thấy có chút là lạ.

\dots

\dots

Cái hành lang này\dots\ sao tôi thấy quen đến vậy?

Sàn gạch màu ngà, trần kẻ ô cùng với một cái loa phát thanh, bên phải là dãy cửa sổ lớn hắt ánh sáng trắng xóa xuống nền, bên trái là những lớp học với bảng thông báo màu xanh bên ngoài dán đầy giấy.

Cảnh này\dots\ tôi đã thấy rồi. Không phải một lần.

\dots Mà là trong giấc mơ hôm qua.

Cảm giác thật đến mức tôi chẳng dám nghĩ là mình mơ nữa.

Nhưng nếu không phải mơ\dots\ thì cái giấc mơ ấy là cái gì? Và tại sao tôi lại bước vào nơi giống hệt nó ngay sáng hôm nay?

Một thoáng rùng mình chạy dọc sống lưng.

Liệu đây có phải là điềm báo? Rằng những gì tôi thấy sẽ sớm trở thành hiện thực?

``\dots Nii-chan! Lớp của chúng ta ở đây này!''

Giọng Saikyou vang lên kéo tôi về thực tại. Tôi giật mình, nhận ra cô bé đã đứng ở cửa lớp từ bao giờ.

``Rồi, anh tới đây,'' tôi đáp, cố giấu đi vẻ bối rối.

Cô em gái của tôi bĩu môi, trách: ``Muộn học tới nơi rồi mà còn lề mề\dots''

Tôi chỉ biết thở dài. Chắc mình nghĩ quá nhiều thôi.

Hy vọng\dots\ ngày đầu tiên ở đây sẽ bình thường như mọi ngôi trường khác.

\begin{center}
  \uline{Lớp 1-C - Trường Trung học Aogami\\Giờ nghỉ giữa giờ.}
\end{center}

Buổi sáng hôm ấy trôi qua với tốc độ nhanh đến mức tôi còn chưa kịp định thần. Tóm gọn lại là thế này\dots

Tiết Ngữ văn đầu tiên với thầy giáo trung niên nghiêm nghị khiến cả lớp im phăng phắc. Tiết Toán thứ hai, cô giáo cho một bài kiểm tra nhỏ để thăm dò trình độ, với Saikyou hăng hái làm bài đến lạ. Tiết Lịch sử thứ ba với thầy giáo trẻ kể chuyện xen bài giảng thú vị, dù tôi vẫn cảm nhận được ánh mắt tò mò của các bạn. Và cuối cùng là tiết Ngoại ngữ sôi động cùng cô giáo trẻ vui tính, khiến em gái tôi thích thú ra mặt mỗi khi trả lời đúng câu hỏi nhanh.

Không khí ở đây dường như cũng khác hẳn so với những ngôi trường trước đó: yên bình và thoải mái hơn hẳn. Không có cảnh chen chúc, xô đẩy như ở các thành phố lớn. Các bạn học cũng có vẻ thân thiện và cởi mở hơn, không phải kiểu ``mạnh ai nấy học'' như những nơi tôi từng trải qua. Ngay cả các thầy cô giáo cũng toát lên vẻ gần gũi, không tạo khoảng cách quá lớn với học sinh như ở những trường khác.

Tiếng chuông giải lao giữa tiết vang lên, một giai điệu nhẹ nhàng nhưng lạ lẫm, khác hẳn kiểu chuông chát chúa đập vào tai ở những ngôi trường trước kia mà hai anh em chúng tôi từng theo học.

Cả lớp bắt đầu ồn ào hơn. Một vài nhóm bạn tụ tập ở cửa sổ, một vài người đi xuống căng-tin, số khác thì ngồi tại bàn trò chuyện. Tôi vừa định hỏi Saikyou có muốn xuống mua gì ăn để lót dạ không, thì bất chợt\dots

``Này, này! Hai cậu cảm thấy thoải mái khi đến Aogami học lần đầu chưa?''

Đó là những câu đầu tiên mà cô gái tóc vàng với hai búi trên đầu nói với chúng tôi ngay khi đã bắt đầu giờ nghỉ. Cô ấy không phải là người duy nhất kéo đến bàn chúng tôi để thăm hỏi - có vẻ như cả lớp đều tò mò về hai học sinh mới - tức chúng tôi.

``Ờm\dots\ cũng chưa hẳn\dots'' tôi đáp, hơi khựng lại trước sự nhiệt tình này.

``Cậu làm anh trai tôi sợ rồi đấy,'' Saikyou xen vào, ánh mắt hơi nheo lại. ``Ờm\dots\ mà bạn tên là gì nhỉ?''

``Thôi nào, Suzune-san,'' một cô gái tóc vàng đằng sau cô ấy cười nhẹ, rồi hướng về chúng tôi với ánh mắt thân thiện, ánh mắt cũng bớt xông xáo: ``Dù gì hôm nay cũng là ngày đầu tiên mà bọn họ tới mà.'' Cô có hai chỏm bím được buộc bằng hai cái vòng hồng viền trắng, nhẹ nhàng đung đưa theo từng cử động.

``Saikyou-san nói đúng đó. Để cho cậu ta có chút thời gian chứ,'' một cô gái tóc xanh lam với bím tết ở bên nhíu mày nhắc khéo. Điều đặc biệt ở cô ấy là một bông hoa trắng cài trên tóc và vài phụ kiện tinh tế ẩn dưới mái tóc.

Cô gái tóc xanh ấy quay sang chúng tôi, khẽ cúi đầu: ``À đúng rồi, bọn mình chưa giới thiệu cho các cậu nhỉ?''

``\dots À phải, giờ nhắc mới nhớ,'' tôi gãi đầu ngượng ngùng. ``Tự dưng các bạn cứ sồn tới thế này làm cho bọn mình thấy khó xử\dots''

Cô gái tóc vàng hai búi cười trừ, cúi rạp đầu xuống như một em bé lễ phép. ``Mình xin lỗi! Vậy để mình giới thiệu trước! Mình là Karamomo Suzune! Rất vui được gặp hai cậu!''

\chrint

Để tôi giới thiệu chi tiết hơn về những người bạn cùng lớp mới của chúng tôi. Họ đều là những gương mặt mới, nhưng\dots\ sao tôi lại có cảm giác quen thuộc đến kỳ lạ?

\begin{wrapfigure}[15]{l}{6.5cm}
  \includegraphics[height=8cm]{../../../../Image/Book?/suzune.png}
\end{wrapfigure}

Karamomo Suzune (\ruby{杏}{からもも}\ruby{鈴}{すず}\ruby{音}{ね}, \ruby{Ca-ra-mô-mô}{An}\ \ruby{Xư-dư}{Linh}-\ruby{ne}{Âm}) là một cô gái tràn đầy năng lượng và nhiệt huyết. Họ của cô được đọc là ``Karamomo'', không phải ``An'' - một cách đọc khá đặc biệt mà ngay cả tôi cũng không hiểu tại sao. Tính cách của cô trong sáng đến mức ngây thơ, luôn háo hức với những điều mới mẻ. Tuy nhiên, đôi khi sự nhiệt tình ấy lại trở thành nhược điểm - hoặc là quá say mê một thứ gì đó, hoặc nhanh chóng chán nản và bỏ cuộc.

Là một người hướng ngoại điển hình, Suzune-san dễ dàng kết bạn với tất cả mọi người. Nhưng như các bạn trong lớp nói, ``nhí nhảnh quá cũng không phải là điều tốt'' - thỉnh thoảng cô ấy gây rắc rối vì không đọc được cảm xúc của mọi người xung quanh.

\chrint

``Ừ, hân hạnh được làm quen,'' em gái tôi đáp lại, bắt tay với Suzune-san.

Tiếp đến là cô gái tóc hai bím, người đang nở nụ cười hiền hậu và thân thiện và đi về phía chúng tôi

``Mình là Ibara Saki. Mình làm việc cho Hội Chăm sóc cây cảnh, và thích chăm sóc các chậu hoa trong trường.''

\chrint

\begin{wrapfigure}[14]{r}{6.5cm}
  \includegraphics[height=8cm]{../../../../Image/Book?/saki_human.png}
\end{wrapfigure}

Ibara Saki-san (\ruby{茨}{いばら}\ruby{咲}{ささ}\ruby{希}{き}, \ruby{I-ba-ra}{Tì}\ \ruby{Xa}{Tiếu}-\ruby{ki}{Hy}) là chị cả trong một gia đình ba chị em kinh doanh hoa cảnh. Hoàn cảnh của cô khá giống chúng tôi - bố mẹ thường xuyên vắng nhà vì công việc, nên cô phải đảm nhận vai trò chăm sóc các em. Có lẽ vì vậy mà Saki-san rất chu đáo và hay giúp đỡ mọi người. Ngoài việc là thành viên Hội Chăm sóc cây cảnh, cô còn giữ chức Ủy viên Hội Văn hóa.

Tình yêu của Saki-san dành cho hoa thể hiện qua cách cô chăm sóc chúng - tỉ mỉ và trân trọng như đối với chính các em của mình. Dưới bàn tay khéo léo của cô, những bông hoa luôn nở rộ rực rỡ hơn bình thường. Quả thực là một người có tâm hồn đẹp và tinh tế.

\chrint

``Còn mình là Yukihara Touka,'' cô gái tóc xanh lam giới thiệu, giọng điềm đạm. ``Chào mừng các cậu đến với lớp học của chúng mình.''

``Ừ, xin mọi người chỉ giáo,'' tôi mỉm cười đáp lại.

\chrint

\begin{wrapfigure}[14]{l}{6.5cm}
  \includegraphics[height=8cm]{../../../../Image/Book?/touka_human.png}
\end{wrapfigure}

Yukihara Touka (\ruby{雪}{ゆき}\ruby{原}{はら}\ruby{冬}{とう}\ruby{花}{か}, \ruby{Giu-ki}{Tuyết}-\ruby{ha-ra}{Nguyện}\ \ruby{Tô-ư}{Đông}-\ruby{ca}{Hoa}) được các bạn trong lớp 1-C công nhận là một thần đồng. Không chỉ xuất sắc trong lớp mà còn nổi tiếng cả khối với trí thông minh nhạy bén. Dù vẻ ngoài có phần lạnh lùng, nhưng bên trong Touka-san lại là một người vô cùng tò mò và ham học hỏi.

Điều thú vị là cô lại rất thân thiết với Suzune-san - một người hoàn toàn đối lập về học lực. Nhưng có lẽ chính sự khác biệt ấy lại khiến tình bạn của họ trở nên đặc biệt hơn. Dù sao thì, tình bạn đâu cần phải dựa trên thành tích học tập?

\chrint

Tôi chợt nhận ra cả ba người họ đang nhìn em gái tôi chăm chú. Saikyou cau mày: ``Bộ mặt tôi có cái gì sao mà mấy người nhìn ghê vậy?''

Ánh mắt của họ\dots\ có điều gì đó không bình thường. Không phải sự tò mò thông thường, mà giống như\dots\ họ đang lo sợ điều gì đó. Tại sao lại như vậy?

\img{e1-1c}

``Nhìn kỹ thì, em gái bạn trông đúng là một người diễm lệ thật,'' Touka-san bất ngờ lên tiếng, nét mặt đã trở lại bình thường. Khả năng chuyển đổi cảm xúc của cô ấy thật đáng nể.

``Diễm lệ\dots?'' Suzune-san nhăn mặt khó hiểu. ``Ể? Nó có nghĩa là gì vậy?''

``Ý của tớ tức là cậu ấy trông xinh đẹp đó.'' Không lạ gì khi Suzune-san không hiểu - đó là một từ cổ hiếm khi được dùng. Đến tôi còn điếng người ra khi nghe cô ấy nói vậy.

``Tôi cũng chả xinh đẹp đến độ để mấy người tâng bốc lắm đâu,'' em gái tôi nhăn mặt, cười nhẹ phản đối. ``Còn thua Nii-chan là đằng khác.''

``Tớ lại thấy hai anh em cậu trông như những búp bê sứ vậy,'' Saki-san nói thêm.

``Bọn mình không tuyệt mỹ đến mức đó\dots'' tôi cười ngượng. Lần đầu tiên chúng tôi nhận được những lời khen như vậy. Môi trường ở đây quả thật khác hẳn những nơi trước đây.

Nhưng\dots\ có phải họ đang quá tốt với chúng tôi không? Nếu nói rằng tôi nghi ngờ thì có vẻ bất công với họ, nhưng cảm giác này thật khó lý giải.

Có lẽ tốt nhất là cứ chấp nhận sự tốt bụng của họ. Đôi khi, im lặng và quan sát mới là lựa chọn khôn ngoan. Dù là những lời xã giao, tôi không thấy họ đang hàm ý khinh miệt gì chúng tôi cả.

``Mà, nhắc đến cái đẹp thì,'' cô con gái bán hoa hướng mắt ra ngoài cửa sổ; chính xác hơn, là cái chậu hoa được đặt ở bên ngoài lớp, ``các cậu có thấy chậu hoa ở gần tủ đựng đồ bên ngoài lớp không?''

\img{e1-1d}

``Ý cậu là chậu hoa đỏ ngoài kia ấy á?'' em gái tôi đáp. ``Nhưng mà, tôi thấy chúng có vẻ quen quen thế nào\dots''

``Ừm\dots\ mình cũng thấy vậy,'' tôi khẽ gật đầu, cố gắng nhớ lại. ``Cứ như đã thấy ở đâu rồi\dots''

Hai anh em chúng tôi cố nhớ lại loài hoa đang được trồng ở ngoài kia. Cô bạn Thần đồng đáp lại cô, đồng thời trả lời cho thắc mắc của chúng tôi: ``Đó là chậu hoa bỉ ngạn đỏ đó. Cũng khá dị đấy.''

``Hoa bỉ ngạn đỏ\dots?'' hai anh em chúng tôi đồng thanh. Trong đầu tôi chợt hiện lên hình ảnh những cánh hoa đỏ thắm, rực rỡ nhưng ma mị trong giấc mơ đêm qua. Không phải một, mà đến tận hai lần - lần đầu là khi chúng tôi bước vào căn phòng có các bàn đặt ở vòng tròn, từ đó đụng độ Shino, và lần thứ hai là khi chúng tôi tìm thấy quyển vở của cô ta.

Cách chọn hoa của Saki-san đúng là dị thật. Thậm chí có phần quái đản. Bởi lẽ\dots

``\dots Mình tưởng đây là loài hoa tượng trưng cho cái chết mà?'' tôi vô thức nói, nhớ lại cảm giác rờn rợn khi nhìn thấy chúng trong mơ.

Chủ nhân của chậu hoa nghe thế, cười nhẹ: ``Đúng như cậu nói. Nhiều người cho rằng chúng là loài hoa khiến họ khó chịu, nhưng không thể phủ nhận rằng chúng rất là đẹp, đúng không?''

``\dots Kể cả vậy, mình không nghĩ là sẽ có ai đó chọn loài hoa này để làm cảnh đâu,'' tôi đáp lại họ với sự lo ngại. Những cánh hoa đỏ thắm trong giấc mơ lại hiện về, cùng với cảm giác lạnh lẽo không thể giải thích.

``Biết đâu sau giờ học, các cậu thử ra đó ngắm kỹ một lần?'' Touka-san nói, nhưng tôi có cảm tưởng rằng đôi mắt của cô ấy trông như người mất hồn. Tiếng cười nhỏ nhẹ lại càng khiến cô ấy trở nên ma mị và khó ngờ hơn.

Không, cả ba cô gái đó đều đang nhìn chúng tôi với cặp mắt vô hồn như thế, như thể có một lớp màn vô hình che đi cảm xúc thật. Sự im lặng giữa câu nói của họ khiến không khí xung quanh nghẹn lại.

Cách hành xử của họ cũng có điều gì đó không ổn. Thậm chí việc chọn loài hoa của con gái chủ cửa hàng bán hoa cũng quá bất thường rồi. Tôi không thể không liên tưởng tới những điều kỳ lạ trong giấc mơ - những bóng người mờ ảo, những ánh mắt trống rỗng, và đặc biệt là những đóa bỉ ngạn đỏ rực.

Dường như suy nghĩ của họ cũng không được bình thường. Tôi cảm tưởng họ đang bị một thế lực nào đó chi phối, giống như cách những bóng ma trong giấc mơ của tôi bị giam cầm trong ngôi trường cũ kỹ ấy.

Tôi cảm thấy họ không giống như là chính mình vậy, dù đây cũng chỉ là lần đầu tôi gặp họ. Nhưng cảm giác này\dots\ quá giống với những gì tôi đã trải qua trong mơ.

Tôi có nghĩ xấu về họ không? Không hề. Nhưng những gì mà tôi nghĩ có phải chỉ là những suy nghĩ vẩn vơ của một kẻ không có bạn không? Cũng không. Đây là linh cảm, một linh cảm mạnh mẽ rằng có điều gì đó không ổn đang diễn ra.

Mọi thứ ở đây đang diễn ra rất bất thường. Đó là điều mà tôi thấy. Nhưng, tôi cũng không muốn đưa ra đánh giá vội. Trước mắt là cứ im lặng và quan sát mọi thứ đã.

Đúng lúc này, tiếng chuông báo hết giờ nghỉ vang lên. Suzune-san nói với chúng tôi bằng giọng nói vui vẻ như chưa hề có chuyện gì xảy ra: ``Vậy, mình hy vọng các bạn sẽ cảm thấy thoải mái hơn trong lớp học nha\~{} Hẹn gặp sau!''

Ba cô gái rời đi, vừa trò chuyện vừa cười khẽ. Nhưng lạ thay, tôi lại có cảm giác như họ\dots\ đang cố tránh nhìn về phía chúng tôi. Sao tôi có cảm tưởng họ đang sợ chúng tôi vậy? Chúng tôi có làm gì sai đâu?

Hay là\dots\ họ cũng cảm nhận được điều gì đó từ chúng tôi, giống như cách chúng tôi cảm nhận được điều bất thường từ họ?

``Này, Nii-chan\dots'' cô em gái nhìn tôi với sự e dè. ``Không biết có phải mỗi mình em không, nhưng\dots''

``\dots Anh biết, anh cũng thấy có gì đó là lạ từ họ,'' tôi nói. Có vẻ như cả Saikyou cũng cảm nhận được sự kỳ lạ từ ba cô gái. ``Giống như trong giấc mơ\dots''

``Em bắt đầu nghi ngờ mấy người học ở đây rồi đó\dots'' em ấy bắt đầu gầm gừ. Tôi ra hiệu bằng cách xua tay lên, cố trấn an: ``Thôi, đừng nóng. Mới có ngày đầu tiên thôi mà. Với cả, ta còn chưa biết mọi người ở nơi này, đúng chứ?''

Em gái tôi dần hạ hỏa, rồi cũng đành miễn cưỡng gật đầu: ``\dots Cũng phải.''

Chúng tôi lặng lẽ lấy sách vở chuẩn bị bài cho tiết tiếp theo. Mong rằng buổi học hôm nay sẽ suôn sẻ\dots

Tôi vẫn có thể nghe được đâu đó Touka-san và Saki-san nói chuyện với nhau về những chậu hoa bỉ ngạn đỏ ấy. Những đóa hoa đỏ thắm, rực rỡ nhưng ma mị, như một lời cảnh báo về điều gì đó sắp xảy ra.

\dots Mà thôi kệ. Dù gì cũng không phải là việc của mình.

Nhưng sâu trong thâm tâm, tôi biết rằng mình không thể phớt lờ nó được. Giấc mơ đêm qua, những đóa bỉ ngạn, và cách cư xử kỳ lạ của ba cô gái\dots\ tất cả đều có liên quan đến nhau, theo một cách nào đó mà tôi chưa thể hiểu được.

\pov{-}

Touka và Saki đi về chỗ ngồi ngay sát cửa sổ của mình, trong khi Suzune thì ngồi ở phía đối diện và làm công việc của mình.

Nàng tóc xanh nghiêng đầu, ánh mắt hướng ra ngoài cửa sổ, giọng nhẹ nhàng: ``Mà này, Saki-chan. Quả thật những bông hoa ngoài kia đẹp thật đấy nhỉ.''

Nàng tóc hai bím lập tức nở nụ cười rạng rỡ, đôi mắt long lanh đầy tự hào.

``Tất nhiên! Chúng rực rỡ đến thế mà!'' cô nói, giọng tràn đầy hãnh diện. ``Chính tay tớ là người đã chăm sóc với tưới nước chúng mỗi ngày đấy.''

Touka hơi mở to mắt, như vừa phát hiện một bí mật nho nhỏ. ``Vậy ra cậu là người chăm sóc chúng à, Saki-san?''

``Chứ còn gì nữa!'' Saki đáp nhanh, vẻ mặt như thể điều này là chuyện hiển nhiên.

Nàng Thần đồng chớp mắt vài lần, rồi khẽ cười. ``Tớ cứ tưởng rằng ai đó ở bên nhân viên bảo vệ làm chứ.''

Saki xua tay lia lịa, mái tóc ngắn khẽ rung theo động tác: ``Không không. Tớ là người làm việc cho Hội Chăm sóc cây cảnh mà, nên tớ là người làm chứ.''

Touka hơi gật gù, nhưng giọng vẫn còn chút ngạc nhiên: ``Bây giờ tớ mới biết cậu làm cho hội ấy đấy. Xin lỗi, tớ không biết.''

``Ơ mà,'' Saki nhướng mày, nghiêng người lại gần, ``tớ tưởng cái lúc hồi người ta phân các hội trong trường, cậu có mặt ở đó chứ, Touka-san? Hay là cậu lại ngủ trong suốt buổi họp hôm đó?''

Cô nàng tóc xanh bật cười nhẹ, lắc đầu: ``Chắc chắn là không rồi. Hồi đó cũng lâu rồi mà, nên tớ cũng quên rồi. Với cả, cậu nghĩ mọi người thực sự sẽ nhớ toàn bộ công việc mà các hội phải làm sao, Saki-chan?''

Saki bỗng im lặng, trầm ngâm vài giây, rồi đáp gọn: ``\dots Không.'' Giọng cô như vừa thừa nhận một điều đơn giản mà mình đã bỏ lỡ.

Touka mỉm cười, hất nhẹ cằm: ``Thấy chưa? Có phải ai cũng nhớ được đâu!''

``Ừ\dots\ Cậu nói cũng phải.'' Saki gật đầu chậm rãi, vẻ mặt hơi lơ đãng.

Nàng tóc xanh bật cười khẽ, giọng trêu chọc: ``Cậu đúng là hơi ngố thật á, Saki-san à.''

Nàng tóc hai bím lập tức nhăn mặt, phản bác ngay: ``Tớ không có ngố!'' rồi như sực nhớ ra điều gì, cô nghiêng đầu, ánh mắt sáng lên ``Mà, Touka-san, cậu có muốn giúp tớ cắt tỉa chậu cây không?''

Touka hơi khựng lại, đôi mày nhíu nhẹ. ``Cắt tỉa á? \dots Hả?''

``Ừ.'' Saki giải thích bằng giọng hào hứng ``Chúng ta phải nhổ bớt vài bông hoa để tránh chúng nó không tranh nhau chất dinh dưỡng chứ. Cứ lấy kéo tỉa và `xoẹt!' mấy cái là xong.'' Cô còn làm động tác cắt bằng tay để minh họa.

Cô bạn đưa tay chạm cằm, ánh mắt hơi dịu xuống. ``Cơ mà như thế cũng hơi tội cho chúng nhỉ.''

Cô nàng chủ cửa hàng hoa khẽ cười, giọng hạ xuống nghe mềm mại hơn: ``Thì thế nên chúng ta sử dụng những bông hoa ấy để cắm vào bình và rồi để chúng vào lớp học chứ.''

Touka nghe vậy thì gật gù, khóe môi nhếch lên thành một nụ cười. ``Cũng ra gì đấy chứ. Được rồi, tớ sẽ giúp.''

``Cảm ơn cậu nha!'' Saki reo lên, giọng đầy phấn khích. ``Sau giờ học hẹn gặp tớ!''

``Đương nhiên rồi,'' Touka đáp, kèm theo một nụ cười nhẹ, như thể câu chuyện nhỏ này vừa khiến ngày học của cô thú vị hơn một chút.

\begin{center}
  \pov{Haragen Saikyou}
  \uline{Một tiếng sau.}
\end{center}

Lại một tiết nữa vừa trôi qua.

Không có gì đặc biệt để kể lại: chúng tôi chỉ đơn giản là ngồi nghe giảng, ghi chép, tiếp thu kiến thức, rồi lặp lại quy trình ấy như một cái máy. Nó khác hẳn với những tiết đầu tiên lúc mới nhập học: khi đó, mọi thứ còn mới mẻ, lạ lẫm, tim tôi đập nhanh hơn một nhịp chỉ vì tiếng bước chân hay lời gọi tên.

Nhưng giờ đây, sự hứng khởi ấy đã nguội lạnh đi nhiều, thay vào đó là cảm giác đơn điệu, như thể đang đi trên một con đường dài mà hai bên chỉ toàn cảnh quen thuộc. Nếu không nhờ học được thói quen dậy sớm từ Nii-chan, có khi tôi đã gục xuống bàn từ nửa tiết trước rồi.

Tiếng chuông vang lên, cả lớp được thả lỏng như một cái lò xo hết căng. Chúng tôi ngồi thư giãn, chờ cho tiết tiếp theo bắt đầu. Giờ giải lao ở đây luôn là thời điểm những nhóm bạn tụ tập, trò chuyện rôm rả. Có nhóm nam nữ lẫn lộn, trêu chọc nhau thoải mái, chẳng cần e dè.

\img{e1-1e}

Còn chúng tôi? Vẫn như mọi khi: ngồi ``tự kỷ'' với nhau. Hoặc, nói nhẹ nhàng hơn, là dành thời gian với nhau theo cách riêng: chơi game điện tử trên điện thoại. Ít nhất, tôi có thể nói rằng chúng tôi đã bắt đầu quen với nhịp sống của môi trường học mới này.

Nhưng rồi, bất ngờ, có hai cô gái tiến về phía bàn chúng tôi. Dáng vẻ của họ trông như muốn bắt chuyện.

Người đầu tiên là một cô tóc nâu sẫm, buộc đuôi ngựa cao, cố định bằng phụ kiện màu đen viền xanh lá, giữa có hình một con mèo trắng. Một chỏm tóc ngố nhỏ vểnh lên trên đỉnh đầu, khiến cô trông thêm phần sinh động.

Đi cạnh cô là một cô gái tóc trắng, buộc đuôi ngựa thấp, bên mái có tết một lọn nhỏ. Trên đầu cũng có chỏm tóc ngố, và cô đội một chiếc bờm có hai tai cáo vểnh nhẹ, một món phụ kiện nổi bật đến mức thoạt nhìn tôi cứ tưởng là tai thật. Nhưng thực tế, cô ấy chỉ có đôi tai người bình thường.

Chúng tôi cùng đặt điện thoại xuống, nhìn họ.

``Chào những người bạn mới chuyển trường!'' cô tóc nâu cất tiếng trước, giọng tươi sáng.

Cô tóc trắng cũng cúi đầu nhẹ, cất lời nhỏ đến mức phải căng tai mới nghe được: ``X-Xin chào\dots''

Tôi hỏi thẳng, chẳng vòng vo: ``Các cậu là ai vậy?'' Chúng tôi thực sự chưa thể nhớ hết tên mặt mọi người trong lớp.

``Tên mình là Hanabi!'' cô tóc nâu trả lời ngay, với sự năng nổ đến mức lan sang cả người nghe. ``Mình là một trong những bạn học mới sẽ giúp đỡ các cậu về mảng học tập!''

\chrint

\begin{wrapfigure}[13]{r}{5cm}
  \includegraphics[height=8cm]{../../../../Image/Book?/hanabi.png}
\end{wrapfigure}

Natsuno Hanabi (\ruby{夏}{なつ}\ruby{乃}{の}\ruby{華}{はな}\ruby{火}{び}, \ruby{Nát-xư}{Hạ}-\ruby{nô}{Nãi}\ \ruby{Ha-na}{Hoa}-\ruby{bi}{Hỏa}) là con gái một gia đình danh giá, và theo lời đồn trong lớp thì cô là người ``truyền lửa'' thật sự. Năng nổ, nhiệt huyết, sẵn sàng giúp đỡ bất cứ ai. Kiểu người đã nói thì sẽ làm, không phải hứa suông. Nhờ có nhiều cơ hội trò chuyện với người lớn trong khu phố, cô ấy sở hữu tính cách mạnh mẽ, tự tin và không hề e ngại bất cứ điều gì.

\chrint

Cô quay sang phía bạn mình, mỉm cười khích lệ: ``Còn đây là Inari! Chào hai cậu ấy đi nào!''

Cô gái tóc trắng - Inari - cúi người lễ phép, nhưng giọng vẫn run và ngập ngừng: ``A-A-nô\dots\ R-Rất vui được gặp hai bạn\dots''

\chrint

\begin{wrapfigure}[13]{l}{4cm}
  \includegraphics[height=8cm]{../../../../Image/Book?/inari_human.png}
\end{wrapfigure}

Shirayuki Inari (\ruby{白}{しら}\ruby{雪}{ゆき}\ruby{稲}{いな}\ruby{裡}{り}, \ruby{Si-ra}{Bạch}-\ruby{giu-ki}{Tuyết}\ \ruby{I-na}{Đạo}-\ruby{ri}{Lý}) thì trái ngược hoàn toàn với cô bạn thuở nhỏ của mình. Lớp trưởng của lớp, nhưng tính tình nhút nhát, như lúc này đây. Hanabi kể rằng cô lớn lên trong một gia đình nghiêm khắc, khiến cô không mấy tự tin vào bản thân, thường dựa vào ý kiến người khác, đặc biệt là Hanabi, cô bạn từ thuở nhỏ.

Việc làm lớp trưởng là quyết định của riêng Inari, với hy vọng thay đổi tính rụt rè ấy\dots\ nhưng thú thật, nếu nói ai hợp vai trò này hơn, tôi sẽ chọn Hanabi-san hơn là cô gái tai cáo rụt rè ấy.

\chrint

``Thôi được rồi, cứ bình tĩnh mà nói. Chúng ta vẫn còn thời gian, nên cứ thoải mái đi,'' Nii-chan mỉm cười, trấn an cô gái tóc trắng kia.

Inari nuốt nước bọt, rồi gồng mình: ``M-Mình có điều này muốn nói\dots''

Bất thình lình, cô bước lên, đập tay xuống bàn *BỘP!* một tiếng khiến cả tôi, Nii-chan và cô bạn giật mình. Tôi thấy rõ cô đang gom hết dũng khí để thốt ra: ``N-Nếu có điều gì muốn hỏi thì cứ bình tĩnh và mạnh dạn hỏi nhé!''

Tôi mỉm cười, vừa đổ mồ hôi vừa đáp: ``Tôi nghĩ câu đó để tôi nói với cậu thì hợp hơn á.''

Nii-chan chỉ cười trừ. ``Không cần phải gắng hết sức đâu. Mình biết bạn cũng đang gắng chí khí của lớp trưởng mà, nhưng không cần phải như thế đâu.''

Hanabi-san nhanh chóng đặt tay lên vai Inari-san, kéo bạn mình về: ``N-Nào, Inari! Bà đang làm họ khó xử đấy.''

Cô bạn tóc trắng giật mình, cúi đầu xin lỗi lần nữa: ``M-M-Mình xin lỗi hai bạn\dots!''

``Không cần cuống thế đâu,''  Nii-chan xua tay, vẫn giữ nụ cười. ``Vẫn còn thời gian giải lao mà. Có gì cứ bình tĩnh nói.''

Cô tóc trắng thở phào, rồi cố nói chuyện bình thường: ``Ờm\dots\ mình là lớp trưởng\dots\ nên là\dots\ ờm\dots\ nếu các cậu có điều gì vướng mắc\dots''

``\dots Bạn nói to hơn được không?'' Nii-chan nhướn mày, giơ tay lên cao ra hiệu.

Tôi còn chưa kịp buông câu ``Lớp trưởng gì mà nói bé thế\dots'', thì Hanabi-san đã nghiêm giọng: ``Nào, Inari! Nói to và rõ ràng hơn nào!''

Cô bắt đầu ``giảng giải'' cho Inari-san như một huấn luyện viên đang thúc đẩy học trò. Trong khi bản thân cô nàng tai cáo thì vừa run rẩy, vừa xanh mặt khi bị đích thân bạn thuở nhỏ chỉ điểm hết lỗi này tới lỗi khác.

Tôi nhìn mà vừa thấy buồn cười, vừa thấy tội cho cô bạn tóc trắng ấy. Với tính cách e dè kia\dots\ thật lòng mà nói, cô ấy chẳng hợp với chức lớp trưởng chút nào.

Sau một hồi ``thuyết giảng'' bạn mình như một giáo viên chủ nhiệm nghiêm khắc, Hanabi-san mới giật mình nhận ra chúng tôi vẫn đang ngồi đó quan sát. Cô khựng lại một giây, rồi nhanh chóng thu lại khí thế hùng hổ vừa rồi, trở lại với nụ cười tươi tắn ban nãy.

``A, xin lỗi hai cậu nha\dots\ Mình mải nói chuyện với bạn ấy quá\dots\ Ehehe\dots'' cô gãi má, cười trừ.

``À, không sao đâu,'' Nii-chan đáp nhẹ nhàng, như thể chẳng có gì. ``Nhìn nét mặt bối rối của Inari-san cũng\dots\ đáng yêu lắm đấy.''

``Đúng không chứ!'' cô bạn năng nổ liền bắt nhịp, hớn hở như thể vừa nhận được đồng minh.

Anh trai ngốc của em thật là\dots\ Anh chỉ giỏi khen con gái người ta thôi. Còn em thì\dots\ chưa bao giờ nghe anh nói em đáng yêu cả. Tôi phồng má lên, chẳng che giấu gì vẻ ghen tị, trong khi Hanabi-san cười tít mắt.

``Cậu ấy thực sự là người có tâm đấy,'' cô bạn nói tiếp, giọng đầy tự hào. ``Chỉ là thỉnh thoảng hơi\dots\ vội vàng quá thôi. Trước đây cậu ấy còn chẳng mở miệng nói câu nào kia!''

Inari khẽ liếc về phía chúng tôi, ánh mắt có phần lúng túng như bị khui ra bí mật thầm kín của mình.

Tôi nghiêng đầu, hỏi thẳng: ``Vậy tại sao cậu lại muốn làm lớp trưởng thế, Inari-san?''

Cô hơi cúi mặt, hai tay đan vào nhau như tìm điểm tựa. Giọng nói run nhẹ, nhưng đủ để nghe rõ: ``M-Mình\dots\ muốn tự tin hơn với bản thân. Nên mình\dots\ đã cố gắng, và\dots''

``\dots tự ứng cử làm lớp trưởng, đúng không?'' tôi tiếp lời. ``Thực ra cũng đáng khen đấy. Nhưng cậu biết mà\dots\ để thay đổi bản thân là một chặng đường dài hơi lắm.''

``Đó, thấy chưa? Ngay cả Saikyou-chan cũng tin là bà làm được!'' Hanabi-san dùng khuỷu tay huých nhẹ vào người bạn thời thơ ấu, đôi mắt lấp lánh sự cổ vũ. ``Tôi tin là bà sẽ làm được thôi!''

``Mình cũng nghĩ vậy,'' Nii-chan góp lời. ``Với Hanabi-san ở bên, mình tin là sẽ có một ngày Inari-san trở thành người mà cả lớp đều quý mến.''

``\dots Cũng\dots\ phải\dots'' Inari đáp, rồi cúi đầu lần nữa. ``Cảm ơn cậu\dots\ nhiều\dots''

Tôi liếc lên đồng hồ treo tường. Vẫn còn khoảng năm phút nữa mới hết giờ giải lao.

``Saikyou, ta vẫn còn chút thời gian,'' Nii-chan quay sang tôi. ``Sao không nhân cơ hội này tìm hiểu nhau nhỉ? Dù sao thì cũng phải làm bạn với ai đó chứ.''

Tôi nhún vai:``Ừ\dots\ Anh nói cũng có lý. Ít ra thì cũng nên mở rộng mối quan hệ.''

Hanabi-san lập tức hưởng ứng: ``Vậy nhé! Các cậu đã biết về bọn mình rồi, giờ đến lượt tụi này tìm hiểu hai cậu chứ ha?''

``Tất nhiên,'' tôi và Nii-chan đồng thanh.

Nhưng Inari-san lại khẽ kéo tay bạn mình, giọng nhỏ nhưng đủ nghe: ``H-Hanabi-chan\dots\ Bà nên bình tĩnh chút\dots''

Tôi bật cười xòa: ``Ôi dào, bình thường thôi mà. Chính cậu nói là nếu có gì muốn hỏi thì cứ hỏi, đúng không? Vậy thì cứ thoải mái.''

Lớp trưởng do dự một chút, rồi khẽ gật đầu: ``\dots Cậu nói cũng đúng.''

Quả thật, việc tìm hiểu lẫn nhau sẽ giúp chúng tôi dễ gắn kết hơn\dots\ về lý thuyết là thế.

Nhưng không hiểu sao, tôi lại thấy có gì đó\dots\ hơi sai sai.

Tại sao họ lại thật sự muốn biết rõ về chúng tôi đến thế?

\dots

\dots

Chắc\dots\ cũng chẳng có gì đáng lo đâu nhỉ? Hai cô bạn mới sau đó kéo ghế và ngồi đối diện chúng tôi, bắt đầu công cuộc tìm hiểu lẫn nhau.

Hanabi-san nghiêng đầu, rồi hỏi với vẻ đầy hứng thú: ``Vậy\dots\ món ăn yêu thích của các cậu là gì?''

Tôi hơi khựng lại. ``Chắc là\dots\ thịt?''

``Cụ thể là thế nào ấy,'' Nii-chan xen vào ngay, giọng đều đều. ``Mình thì thích thịt bò.''

Tôi nhún vai. ``Tôi thì\dots\ ăn gì cũng được.''

Hanabi-san gật gù như ghi chú lại điều gì đó trong đầu. ``Vậy còn môn thể thao yêu thích?''

Anh tôi đưa tay gãi má, có chút bối rối. ``Mình không hay chơi thể thao lắm\dots\ Nhưng mà nếu chọn thì chắc là\dots\ chạy bộ?''

Tôi bật cười khẩy, như tìm được cơ hội để cà khịa: ``Anh lười chảy ì ra thì làm gì có môn thể thao nào anh thích được chứ.''

``Em biết anh yếu thể lực thế nào rồi mà\dots'' anh vừa nói vừa đảo mắt sang hướng khác, rõ là không muốn tranh luận thêm.

``Cậu thì sao, Saikyou-chan?'' Hanabi-san quay sang tôi.

``Mình thì chơi tất. Môn gì cũng được.'' Tôi trả lời ngắn gọn, không mấy để tâm.

Inari-san vẫn im lặng nãy giờ, ánh mắt chỉ đảo qua chúng tôi rồi lại cúi xuống bàn. Nhưng tôi thấy được ánh mắt ngưỡng mộ của cô ấy, có lẽ vậy.

``À mà\dots'' cô bạn năng nổ bỗng nghiêng người về phía trước, giọng rộn ràng. ``Trò chơi điện tử yêu thích của các cậu là gì?''

``Âm nhạc và giải đố,'' tôi và Nii-chan đồng thanh đáp.

``Ha! Hai người họ đồng thuận luôn kìa!'' Inari-san bật thốt, dù giọng vẫn nhỏ nhẹ.

Hanabi-san cười khoái chí, rồi hỏi tiếp: ``Điểm nào mà các cậu thấy thu hút nhất?''

``Ờm\dots'' anh tôi ngập ngừng, liếc sang tôi tìm cứu viện.

``Cái này thì\dots\ hơi khó nói ha,'' tôi đáp, thành thật.

``Với cả,'' anh nói thêm, ``bọn mình cũng thấy cái này hơi không liên quan nhỉ\dots''

Cô Lớp trưởng bỗng khẽ kéo tay áo cô bạn của mình, như muốn phản kháng: ``H-Hanabi-chan à\dots''

``Sao thế?'' Hanabi-san quay đầu lại, vẻ như chưa hiểu. ``Chẳng phải chúng ta cần phải biết về họ nhiều nhất có thể à?''

``Tôi thì không ngại đâu,'' tôi chen vào, ``nhưng tại sao cậu lại cần nhiều thông tin từ chúng tôi vậy?''

``Để chuẩn bị cho bữa tiệc chào mừng vào cuối tuần sau á!'' Hanabi-san đáp ngay, như thể đó là điều hiển nhiên.

Inari-san khẽ gật đầu, rồi nói nhỏ: ``Thì, đúng là thế, nhưng mà\dots\ Chúng ta đâu có cần phải biết điểm gì thu hút từ họ đâu.''

``Mình cũng thấy cái đó không liên quan lắm tới việc chuẩn bị bữa tiệc ấy,'' Nii-chan phụ họa.

``Tôi thậm chí còn có cảm giác là cậu đang định rình mò chúng tôi á,'' tôi nói nửa đùa nửa thật.

Cô bạn năng nổ đưa tay lên làm động tác đầu hàng:  ``Quái\dots\ Bắt quả tang luôn cả tôi rồi kìa. Này, tôi thực sự tò mò về họ thật mà! Chỉ là làm quá thôi, làm gì căng!''

``Mình nghĩ chắc bạn cũng đang khá tò mò, đúng chứ?'' anh tôi hỏi, ánh mắt hướng sang Inari-san.

``À-À thì, mình k-không\dots'' cô nàng tóc trắng đỏ mặt, lảng tránh ánh nhìn từ phía chúng tôi.

``Tôi hỏi thật này,'' tôi nghiêng đầu, ``tại sao cậu lại hỏi cặn kẽ chúng tôi đến vậy?''

Hanabi-san chớp mắt một cái, rồi nở nụ cười tươi kèm một cú nháy mắt đầy ẩn ý: ``Nói thật thì, mình luôn bị thu hút bởi những người tốt bụng, sẵn sàng giúp đỡ người ta khi họ cần.''

``Kiểu như Nii-chan á?'' tôi bật miệng. ``Anh ấy cũng là người như vậy mà.''

``Chắc thế-- Này, sao lại lôi anh vào!?'' anh tôi phản ứng ngay, hơi đỏ mặt.

``Thôi, thôi, Hanabi-chan,'' Inari-san kéo tay bạn mình lại, ``giờ không phải lúc để mộng tưởng đâu.''

``Trời ạ, Inari, bà cứ nghiêm trọng hóa mọi thứ thế\dots'' Hanabi-san cười khúc khích rồi áp sát, hai tay đặt lên vai bạn mình, bắt đầu đung đưa. ``Thả lỏng suy nghĩ của mình đi chứ! Thả lỏng đi nào!''

``Ha-na-bi-chan à\dots!'' Lớp trưởng rên khẽ, rõ ràng là đang hơi choáng.

Nii-chan bật cười. ``Hai cậu trông có vẻ thân thiết thật đấy.''

``Cơ mà nhẹ tay thôi, kẻo Inari-san chóng mặt đấy,'' tôi nhắc.

``Ối chà, xin lỗi hai cậu nha.'' Hanabi-san rụt tay lại, rồi cười hồn nhiên. ``Được rồi, câu hỏi tiếp theo--''

*REEEEENG!*

Tiếng chuông báo hiệu vang lên, cắt ngang câu nói.

``Ồ, hết giờ nghỉ rồi kìa,'' Nii-chan nói, mắt ngước lên đồng hồ.

``Nào, Hanabi-chan. Thầy vào bây giờ đấy,'' Inari-san kéo bạn mình về chỗ.

``Hay là sau giờ học chúng ta nói chuyện tiếp nhỉ?'' tôi đề nghị.

``Ý hay á!'' Hanabi-san lập tức đồng ý. ``Lúc đó mình sẽ tiếp tục hỏi các cậu nhiều câu hơn nữa, nên là các cậu chuẩn bị tinh thần đi nhá!''

``Miễn là các cậu đừng hỏi cái gì khiếm nhã quá là được,'' tôi đáp, nửa đùa, nửa thật.

``V-Vậy thì một lần nữa\dots'' Inari-san cúi đầu, giọng chậm rãi nhưng rõ ràng. ``Rất vui được gặp các cậu tại đây. Mình muốn các cậu hãy mong chờ tới bữa tiệc chào mừng sắp tới.''

``Ừ, mình cũng hy vọng là như vậy,'' anh tôi đáp, khẽ cúi đầu xuống như lấy lệ.

``Đừng để cho bọn này thất vọng đấy,'' tôi thêm vào.

Và cứ thế, hai cô bạn ấy vui vẻ trở lại chỗ ngồi. Niềm nở đến mức tôi cũng phải công nhận là\dots\ hiếu khách thật.

\dots Nhưng mà, dù nụ cười trên môi họ trông rạng rỡ đến đâu, khi nhìn kỹ vào đôi mắt ấy, tôi lại thấy một điều khác: một tia cầu cứu mơ hồ, như thể họ đang phải đóng vai này vì một lý do nào đó. Cả cuộc trò chuyện vừa rồi\dots\ có gì đó hơi giả.

Không đời nào họ lại tốt với chúng tôi vô điều kiện như vậy. Hẳn là có nguyên do. Không, tôi không nghĩ đó là dụng ý xấu xa. Chỉ là\dots\ có cảm giác một thế lực nào đó đang buộc họ phải làm vậy.

Nii-chan nói đúng. Nơi này\dots\ không bình thường. Nhưng, tôi không muốn làm khó họ. Ít nhất, họ vẫn cho thấy thành ý.

Có lẽ, chúng tôi cứ nên quan sát thêm một thời gian nữa.

Được rồi, chuẩn bị cho tiết tiếp theo thôi.

\pov{-}

``Ufufu\~{}\dots''

Dưới hành lang tĩnh lặng, bên ngoài cánh cửa lớp học, một tiếng cười khẽ vang lên, nghe như thể đang ẩn giấu niềm vui khó kìm nén.

Đứng đó là một cô gái với mái tóc nâu dài đến giữa lưng, từng lọn tóc khẽ lay động theo làn gió len qua khung cửa sổ. Đôi mắt nâu sáng long lanh của cô hướng về phía trong lớp, nơi cặp sinh đôi mới chuyển trường đang ngồi giữa những người bạn mới.

Ánh nhìn của cô không phải là sự tò mò hời hợt; nó sâu và đầy ý nghĩa, như thể từng giây từng phút quan sát đều đã được chờ đợi từ rất lâu. Môi cô khẽ cong lên, nụ cười lan rộng, mang một thứ ấm áp hiếm có, một niềm hạnh phúc thuần khiết, không chút giả tạo.

``Cuối cùng thì\dots'' cô thì thầm, quay đầu nhìn về khoảng hành lang dài thẳng tắp, nắm chặt tay mình như để kìm lại dòng cảm xúc đang dâng trào. Trái tim cô đập nhanh hơn, từng nhịp như đang đếm ngược tới khoảnh khắc quan trọng nhất đời.

``Ước nguyện\dots\ bấy lâu nay\dots'' cô khép mắt lại một giây, để mặc cho hơi thở hòa cùng nhịp đập hân hoan. Rồi đôi mắt nâu ấy mở ra, ánh sáng kiên định bừng lên.

``Mình đã có thể thực hiện nó rồi. Ước nguyện mà mình vẫn giữ khư khư trong tim\dots\ cuối cùng cũng đến lúc!''

Ở phía xa, tiếng chuông báo hiệu tiết học mới vừa vang lên. Nhưng với cô, mọi âm thanh xung quanh đều như mờ nhạt. Tất cả chỉ còn lại cảm giác, rằng thời khắc mà cô chờ đợi đã thật sự tới.

\end{document}